% ==============================================================================
% CAPÍTULO 2 - METODOLOGIA
% ==============================================================================

\section{Metodologia}
\label{sec:metodologia}

Esta pesquisa é de natureza aplicada, com abordagem quantitativa e procedimento experimental:
um sistema é construído e avaliado em laboratório, com dados gerados de forma controlada e
métricas definidas antes da análise. Quanto aos objetivos, é exploratória, ao investigar se o
log nativo basta para detectar acessos anômalos, e descritiva, ao medir com que precisão e
cobertura isso acontece. Uma pesquisa bibliográfica sustenta a fundamentação e a comparação
com os trabalhos relacionados.

A escolha do experimento em laboratório decorre de uma restrição prática: não existem bases
públicas rotuladas de acessos a bancos de dados com anomalias conhecidas, e rótulos confiáveis
são indispensáveis para medir o desempenho de um detector. Por isso, todo acesso avaliado é
gerado por um simulador, executado de verdade em um servidor PostgreSQL e registrado pelo
próprio servidor, e os rótulos das anomalias ficam em um manifesto à parte, que o sistema
avaliado nunca lê.

O trabalho seguiu sete etapas, cada uma concluída com uma revisão antes da seguinte:

\begin{enumerate}[label=\alph*)]
    \item \textbf{revisão bibliográfica} sobre segurança de bancos de dados, o registro de log
          do PostgreSQL, detecção de anomalias e avaliação de detectores;
    \item \textbf{laboratório e instrumentação do log}: servidor em contêineres, base de dados
          de exemplo, configuração do registro de log e medição do seu custo;
    \item \textbf{simulação}: usuários com rotinas e ruído legítimo, anomalias injetadas em
          várias intensidades e verificações automáticas de coerência entre o que foi executado
          e o que foi registrado;
    \item \textbf{desenvolvimento do SACLA}: leitura do log, extração de eventos, atributos por
          janela, modelo, limiar, explicação dos alertas e relatório;
    \item \textbf{avaliação experimental}: divisão temporal em treino, validação e teste;
          ajuste de hiperparâmetros só na validação; teste usado uma vez; comparação com regras
          de especialista e com a combinação das duas abordagens;
    \item \textbf{verificação da solidez}: réplicas com outras sementes de simulação; uma
          validação confirmatória em simulações novas, com hipóteses registradas antes dos
          dados; uma ablação dos grupos de atributos; e um cenário com mudanças legítimas na
          rotina dos usuários, também pré-registrado;
    \item \textbf{documentação e reprodutibilidade}: cada decisão técnica foi registrada com a
          sua fonte e evidência, e todo o experimento é reproduzível por comandos.
\end{enumerate}

As métricas são as usuais para detectores binários com classes desbalanceadas: precisão,
recall e F1 por janela de tempo, área sob as curvas ROC e precisão--recall, e a fração de
cenários anômalos detectados, em que um cenário conta como detectado se ao menos uma de suas
janelas gera alerta. Para as comparações entre detectores, usam-se intervalos de confiança de
95\% por \textit{bootstrap}, testes de McNemar e de permutação pareada e a correção de Holm para
comparações múltiplas, descritos no Referencial Teórico.

O pré-registro consistiu em fixar, num \textit{commit} datado do repositório do projeto e antes
de gerar qualquer dado novo, a configuração congelada do modelo, as sementes das simulações, as
hipóteses, os testes e o nível de significância, com o compromisso de relatar os resultados
como saíssem \cite{nosek2018}. Todo o processamento roda em contêineres com versões fixadas, as
sementes aleatórias são fixas e o código tem 206 testes automatizados; a cada etapa, os números
das etapas anteriores foram reproduzidos exatamente, com diferença máxima da ordem de
$10^{-16}$. Um assistente de IA generativa foi usado como apoio na implementação, na execução
dos comandos e na redação, conforme a declaração ao final do trabalho; os resultados vêm da
execução do código e podem ser reproduzidos pelos comandos do Apêndice B.
